
\documentclass{article}
\usepackage{colortbl}
\usepackage{makecell}
\usepackage{multirow}
\usepackage{supertabular}

\begin{document}

\newcounter{utterance}

\centering \large Interaction Transcript for game `cladder', experiment `full\_v1.5\_default', episode 6572 with Qwen3.5{-}27B.
\vspace{24pt}

{ \footnotesize  \setcounter{utterance}{1}
\setlength{\tabcolsep}{0pt}
\begin{supertabular}{c@{$\;$}|p{.15\linewidth}@{}p{.15\linewidth}p{.15\linewidth}p{.15\linewidth}p{.15\linewidth}p{.15\linewidth}}
    \# & \multicolumn{2}{c}{Player} && \multicolumn{2}{c}{Game Master} \\
    \hline

    \theutterance \stepcounter{utterance}  
    & & & \multicolumn{4}{p{0.6\linewidth}}{
        \cellcolor[rgb]{0.9,0.9,0.9}{
            \makecell[{{p{\linewidth}}}]{
                \texttt{\tiny{[P1$\langle$GM]}}
                \texttt{You are an expert in causal inference. The following question is not a typical commonsense query, but rather a meticulously designed question created by a professor specializing in causal inference, intended to assess the students' mastery of the course content.} \\
\texttt{Imagine a self{-}contained, hypothetical world with only the following conditions, and without any unmentioned factors or causal relationships: Kwox has a direct effect on kwoz. Swoq has a direct effect on kwoz. For those who are not kwox, the probability of kwoz is 82\%. For those who are kwox, the probability of kwoz is 32\%. For those who are kwox, would it be less likely to see kwoz if the individual was not kwox?} \\
\texttt{Start your answer with 'Yes' or 'No', followed by additional reasoning or evidence to support your explanation.} \\
            }
        }
    }
     \\ \\

    \theutterance \stepcounter{utterance}  
    & \multicolumn{4}{p{0.6\linewidth}}{
        \cellcolor[rgb]{0.561,0.737,0.561}{
            \makecell[{{p{\linewidth}}}]{
                \texttt{\tiny{[P1$\rangle$GM]}}
                \texttt{Yes} \\
\\ 
\texttt{The question asks whether, for an individual who is currently **kwox**, the probability of observing **kwoz** would be lower if that individual were instead **not kwox**.} \\
\\ 
\texttt{Based on the provided conditions:} \\
\texttt{1.  The probability of **kwoz** given **kwox** is **32\%** (\$P(\textbackslash{}text\{kwoz\} | \textbackslash{}text\{kwox\}) = 0.32\$).} \\
\texttt{2.  The probability of **kwoz** given **not kwox** is **82\%** (\$P(\textbackslash{}text\{kwoz\} | \textbackslash{}neg\textbackslash{}text\{kwox\}) = 0.82\$).} \\
\\ 
\texttt{The question essentially compares the counterfactual state (being not kwox) to the factual state (being kwox) for the same individual. Since the probability of **kwoz** drops from 82\% to 32\% when the condition changes from "not kwox" to "kwox," it follows that if the individual were **not kwox**, the likelihood of seeing **kwoz** would be **higher** (82\%) than it is in their current state (32\%).} \\
\\ 
\texttt{However, the phrasing "would it be **less likely** to see kwoz if the individual was not kwox?" requires checking the direction of the inequality.} \\
\texttt{{-} Current state (kwox): 32\% chance.} \\
\texttt{{-} Hypothetical state (not kwox): 82\% chance.} \\
\\ 
\texttt{Is 82\% **less likely** than 32\%? No, 82\% is greater than 32\%. Therefore, if the individual were not kwox, it would be **more** likely to see kwoz, not less.} \\
\\ 
\texttt{Wait, let me re{-}read the specific phrasing carefully: "For those who are kwox, would it be less likely to see kwoz if the individual was not kwox?"} \\
\\ 
\texttt{This is a counterfactual comparison:} \\
\texttt{{-} \$P(\textbackslash{}text\{kwoz\} | \textbackslash{}text\{do\}(\textbackslash{}neg\textbackslash{}text\{kwox\})) = 0.82\$} \\
\texttt{{-} \$P(\textbackslash{}text\{kwoz\} | \textbackslash{}text\{do\}(\textbackslash{}text\{kwox\})) = 0.32\$} \\
\\ 
\texttt{The question asks: Is \$P(\textbackslash{}text\{kwoz\} | \textbackslash{}neg\textbackslash{}text\{kwox\}) < P(\textbackslash{}text\{kwoz\} | \textbackslash{}text\{kwox\})\$?} \\
\texttt{Is \$0.82 < 0.32\$?} \\
\texttt{No.} \\
\\ 
\texttt{Therefore, the answer is **No**. It would be **more** likely to see kwoz if the individual was not kwox.} \\
\\ 
\texttt{**Correction on the initial "Yes" trigger:**} \\
\texttt{The prompt asks me to start with "Yes" or} \\
            }
        }
    }
    & & \\ \\

    \theutterance \stepcounter{utterance}  
    & & & \multicolumn{2}{p{0.3\linewidth}}{
        \cellcolor[rgb]{0.9,0.9,0.9}{
            \makecell[{{p{\linewidth}}}]{
                \texttt{\tiny{[GM$|$GM]}}
                \texttt{Yes} \\
            }
        }
    }
    & & \\ \\

    \theutterance \stepcounter{utterance}  
    & & & \multicolumn{2}{p{0.3\linewidth}}{
        \cellcolor[rgb]{0.9,0.9,0.9}{
            \makecell[{{p{\linewidth}}}]{
                \texttt{\tiny{[GM$|$GM]}}
                \texttt{no} \\
            }
        }
    }
    & & \\ \\

    \theutterance \stepcounter{utterance}  
    & & & \multicolumn{2}{p{0.3\linewidth}}{
        \cellcolor[rgb]{0.9,0.9,0.9}{
            \makecell[{{p{\linewidth}}}]{
                \texttt{\tiny{[GM$|$GM]}}
                \texttt{game\_result = LOSE} \\
            }
        }
    }
    & & \\ \\

\end{supertabular}
}

\end{document}
